\chapter{AI-DLC: how Slotty is changed}\label{ch:aidlc}
\chapterstatus{complete}{2026-11-16}{\src{.claude}, \src{aidlc}, AI-DLC Workflows 2.11.4}

Slotty is changed with the \term{AI-Driven Development Life Cycle} (AI-DLC), a method from AWS
in which an AI agent proposes and a person decides. For every activity the agent writes a plan,
asks the questions it needs answered, and does the work only after a person approves; every
question, answer and approval is kept in an audit trail inside the repository. This chapter
explains the method as installed here, how to run a change with it, and which chapters of this
document each part of it reads and updates.

\section{Vocabulary}
\begin{description}
  \item[Intent] One piece of work, stated as the outcome wanted (\enquote{students can swap slots
    with each other}). Each intent gets its own record and audit trail.
  \item[Phase and stage] The method's five phases contain 33 stages. Some always run; the rest run
    only when the work needs them (\cref{tab:aidlc-stages}).
  \item[Unit of work] A piece of the intent that can be built and tested on its own, the AI-DLC
    counterpart of an epic.
  \item[Bolt] A short build cycle, hours or days rather than weeks, the counterpart of a sprint.
  \item[Gate] A point where the workflow stops for a person's approval, with the evidence to
    decide on.
  \item[Profile (scope)] A preset choosing which stages run and how deep, from a five-stage
    proof of concept to the full thirty-stage lifecycle (\cref{tab:aidlc-profiles}).
\end{description}

\section{The lifecycle}
\begin{figure}[htbp]
  \centering
  % AI-DLC Workflows 2.11 lifecycle (chapter 13): five phases, their stages, and the gates.
\begin{tikzpicture}[x=31.5mm, y=1mm,
    stg/.style={font=\sffamily\scriptsize, text width=29mm, align=left, anchor=north, inner sep=1pt, execute at begin node={\hyphenpenalty=10000}}]
  \foreach \i/\name/\tone/\n in {0/Initialization/notetone/3, 1/Ideation/todotone/7,
      2/Inception/decisiontone/9, 3/Construction/ruletone/7, 4/Operation/audittone/7} {
    \node[phase=\tone, text width=27mm] (p\i) at (\i,0) {\name\\[-1pt]{\normalfont\tiny\color{muted}\n\ stages}};
  }
  \foreach \i in {0,...,3} \pgfmathtruncatemacro{\j}{\i+1} \draw[flow] (p\i) -- (p\j);
  \node[stg] at (0,-7) {Workspace scaffold\\Workspace detection\\State initialization\\[2pt]\textit{no gates}};
  \node[stg] at (1,-7) {Intent capture\\Market research\\Feasibility\\Scope definition\\Team formation\\Rough mockups\\Approval \& handoff};
  \node[stg] at (2,-7) {Reverse engineering\\Practices discovery\\Requirements analysis\\User stories\\Refined mockups\\Domain design\\Units generation\\Contract design\\Delivery planning};
  \node[stg] at (3,-7) {Functional design\\NFR requirements\\NFR design\\Infrastructure design\\Code generation\\Build and test\\CI pipeline};
  \node[stg] at (4,-7) {Deployment pipeline\\Environment provisioning\\Deployment execution\\Observability setup\\Incident response\\Performance validation\\Feedback \& optimization};
  \node[font=\sffamily\scriptsize, text=muted, align=center] at (2,-52) {AI proposes a plan and asks $\rightarrow$ a person answers and approves at each gate $\rightarrow$ AI executes;\\every question, answer and approval is written to the intent's audit trail};
\end{tikzpicture}

  \caption{AI-DLC's five phases. Each stage ends at a gate unless it is a bootstrap stage; the
    agent carries every approved artifact forward as context for the next.}
  \label{fig:aidlc}
\end{figure}

\begin{description}
  \item[Initialization] Creates the intent's record, classifies the workspace (Slotty is
    \emph{brownfield}) and sets up its state. Runs without gates.
  \item[Ideation] Frames the intent, its feasibility and its scope, and ends with an initiative
    brief for approval. Most Slotty work skips it (\code{classic} profile).
  \item[Inception] Turns the intent into requirements, user stories, a domain design and units
    of work, then plans delivery. On Slotty, \emph{Reverse Engineering} builds a code knowledge
    base first, and \emph{Practices Discovery} confirms the project's rules.
  \item[Construction] Designs each unit (functional and non-functional), generates the code and
    tests, builds and tests everything, and updates CI when needed.
  \item[Operation] Deployment pipeline, environments, deployment, observability, incident
    response, performance validation and feedback.
\end{description}

{\tablefont
\begin{xltabular}{\linewidth}{@{}l L l@{}}
  \caption{The 33 stages of AI-DLC Workflows 2.11.4, from the compiled stage graph
    (\texttt{.claude/tools/data/stage-graph.json}).}\label{tab:aidlc-stages}\\
  \toprule
  \thead{No.} & \thead{Stage} & \thead{Runs}\\
  \midrule
  \endfirsthead
  \toprule
  \thead{No.} & \thead{Stage} & \thead{Runs}\\
  \midrule
  \endhead
  0.1 & Workspace Scaffold & always\\
  0.2 & Workspace Detection & always\\
  0.3 & State Initialization & always\\
  1.1 & Intent Capture \& Framing & always\\
  1.2 & Market Research & when needed\\
  1.3 & Feasibility \& Constraints & when needed\\
  1.4 & Scope Definition & always\\
  1.5 & Team Formation & when needed\\
  1.6 & Rough Mockups & when needed\\
  1.7 & Approval \& Handoff & always\\
  2.1 & Reverse Engineering & when needed\\
  2.2 & Practices Discovery & always\\
  2.3 & Requirements Analysis & always\\
  2.4 & User Stories & when needed\\
  2.5 & Refined Mockups & when needed\\
  2.6 & Domain Design & when needed\\
  2.7 & Units Generation & always\\
  2.8 & Contract Design & when needed\\
  2.9 & Delivery Planning & always\\
  3.1 & Functional Design & when needed\\
  3.2 & NFR Requirements & when needed\\
  3.3 & NFR Design & when needed\\
  3.4 & Infrastructure Design & when needed\\
  3.5 & Code Generation & always\\
  3.6 & Build and Test & always\\
  3.7 & CI Pipeline & when needed\\
  4.1 & Deployment Pipeline & when needed\\
  4.2 & Environment Provisioning & when needed\\
  4.3 & Deployment Execution & when needed\\
  4.4 & Observability Setup & when needed\\
  4.5 & Incident Response & when needed\\
  4.6 & Performance Validation & when needed\\
  4.7 & Feedback \& Optimization & when needed\\
  \bottomrule
\end{xltabular}}

\section{Profiles}
{\tablefont
\begin{xltabular}{\linewidth}{@{}l r L L@{}}
  \caption{Workflow profiles, and when to use each on Slotty.}\label{tab:aidlc-profiles}\\
  \toprule
  \thead{Profile} & \thead{Stages} & \thead{What it runs} & \thead{Use on Slotty for}\\
  \midrule
  \endfirsthead
  \toprule
  \thead{Profile} & \thead{Stages} & \thead{What it runs} & \thead{Use on Slotty for}\\
  \midrule
  \endhead
  \code{classic} & 15 & Inception and Construction, ending at Build and Test (default) & Most feature work\\
  \code{feature} & 30 & The full lifecycle, Ideation to Operation & A large feature that changes infrastructure too\\
  \code{bugfix} & 6 & Understand, requirements, code, build and test, deploy & A defect\\
  \code{refactor} & 7 & A behaviour-preserving internal change & Restructuring without changing rules\\
  \code{infra} & 12 & Requirements, NFRs, infrastructure design, CI/CD, deployment, observability & Changes under \src{infra/aws}\\
  \code{security-patch} & 7 & Security requirements through deployment & A vulnerable dependency, a security finding\\
  \code{express} & 7 & Requirements to code and tests, minimal depth & Small, low-risk changes\\
  \code{poc} & 5 & A feasibility check & A spike before committing to an approach\\
  \code{mvp} & 20 & Full Inception and Construction, no Operation & A new product area\\
  \code{enterprise} & 30 & Full lifecycle, comprehensive depth and tests & Regulated or high-assurance work\\
  \code{workshop} & 23 & Facilitated sessions, one advisory review & Team sessions (mob elaboration)\\
  \bottomrule
\end{xltabular}}

\section{Running a change}
Start Claude Code in the repository and describe the outcome:
\begin{lstlisting}
/aidlc students can swap their booked slots with each other
/aidlc --doctor                      # check the installation
/aidlc compose "fix the reminder time in the course timezone"   # propose a tailored plan first
/aidlc --stage domain-design --single  # run one stage on its own
\end{lstlisting}
The agent picks a profile (or proposes one with \code{compose}), and works stage by stage. At each
gate it shows what it produced and waits. Construction asks once whether to \ui{Continue
automatically} or \ui{Review each checkpoint}, and always asks for approval of each unit's code plan.

\subsection{Where things are kept}
\begin{xltabular}{\linewidth}{@{}P{58mm} L@{}}
  \toprule
  \thead{Path} & \thead{Contents}\\
  \midrule
  \code{aidlc/spaces/default/intents/<YYMMDD>-<label>/} & One intent: \code{aidlc-state.md} (where it is), the audit trail, and every stage's artifacts under \code{<phase>/<stage>/}\\
  \code{aidlc/spaces/default/codekb/} & The code knowledge base from Reverse Engineering, reused while it is fresh\\
  \code{aidlc/spaces/default/memory/project.md} & Slotty's standing rules for the agent: stack, architecture rules, decisions, mandated and forbidden practices\\
  \code{aidlc/spaces/default/memory/phases/} & Per-phase guidance\\
  \src{.claude} & The installed framework: orchestrator skill, 14 agent roles, hooks, sensors, tools. Generated by \code{aidlc config}; do not edit by hand\\
  \bottomrule
\end{xltabular}
Intent records are committed with the change they describe, so the reasoning travels with the
code. Local working files (the active intent, sessions, diagnostics) are git-ignored.

\begin{auditnote}
  For any change made with AI-DLC, the intent folder shows the requirement as approved, the
  questions asked and their answers, each gate's approval, and the plan the code followed. Compare
  its \emph{Delivery Planning} artifact with the pull request.
\end{auditnote}

\section{How AI-DLC keeps this document current}
Each phase reads the chapters that describe what it is about to change, and its gates require the
same chapters to be updated (\cref{fig:aidlc-docs}). This is written into the project rules the
agent loads (\code{aidlc/spaces/default/memory/project.md}), so it applies to every intent.

\begin{figure}[htbp]
  \centering
  % Phase × chapter matrix (chapter 13): r = reads, u = must update when affected.
\begin{tikzpicture}[x=17mm, y=-5.2mm, font=\sffamily\scriptsize]
  \foreach \i/\p in {1/Init., 2/Ideation, 3/Inception, 4/Construction, 5/Operation}
    \node[font=\sffamily\scriptsize\bfseries, anchor=base] at (\i,-0.1) {\p};
  \foreach \r/\ch/\cells in {
      1/{1 Introduction}/{r,,,,},
      2/{2 Context, quality goals}/{,u,r,u,r},
      3/{3 Roles, journeys}/{,r,u,r,},
      4/{4 Business rules}/{,,u,u,},
      5/{5 Building blocks}/{,,r,u,},
      6/{6 Data model}/{,,r,u,},
      7/{7 Cross-cutting}/{,,r,u,},
      8/{8 AWS platform}/{,,,u,u},
      9/{9 Delivery}/{,,,u,u},
      10/{10 Operations}/{,,,r,u},
      11/{11 Security}/{,,r,u,r},
      12/{12 Cost}/{,r,,u,u},
      13/{13 AI-DLC}/{r,,,,},
      14/{14 Decisions}/{,r,u,u,r},
      15/{15 Roadmap}/{,u,r,,u}} {
    \node[anchor=east] at (0.45,\r) {\ch};
    \foreach \c [count=\i] in \cells {
      \ifx\c\empty\else
        \if\c u \fill[accent] (\i,\r) circle (1.5mm);\else \draw[accent, line width=.6pt] (\i,\r) circle (1.4mm);\fi
      \fi
    }
    \ifodd\r \begin{scope}[on background layer]\fill[wash] (-2.2,\r-.5) rectangle (5.5,\r+.5);\end{scope}\fi
  }
\end{tikzpicture}

  \caption{Which chapters each AI-DLC phase reads (hollow dot) and must update when the change affects them (filled dot).}
  \label{fig:aidlc-docs}
\end{figure}

\section{Installation}
\begin{enumerate}
  \item Install the \code{aidlc} command once per machine. The installer verifies SHA-256 release
    checksums, installs into \code{~/.local/bin} and \code{~/.local/share/aidlc}, and needs no
    administrator rights:
\begin{lstlisting}[language=bash]
curl -fsSL https://github.com/awslabs/aidlc-workflows/releases/latest/download/install.sh -o install.sh
less install.sh && sh install.sh
\end{lstlisting}
  \item The repository is already configured for Claude Code (\code{aidlc config --harness claude});
    \code{aidlc doctor} checks it. Start Claude Code from a terminal so its hooks find
    \code{aidlc} on the PATH; if you start it from a desktop launcher, add \code{~/.local/bin} to
    the PATH in \code{~/.config/environment.d/}.
  \item Approve the project hooks when Claude Code asks, then restart it.
\end{enumerate}
